\documentclass[aspectratio=169,12pt]{beamer}
\usepackage{amsmath}
\usepackage{amssymb}
\usepackage{array}
\usepackage[most]{tcolorbox}
\usepackage{graphicx}
\usepackage[sfdefault,lf]{FiraSans}
\renewcommand{\familydefault}{\sfdefault}
\setlength{\parindent}{0pt}
\everymath{\displaystyle}
\setbeamertemplate{navigation symbols}{}
\setbeamertemplate{footline}{}
\setbeamertemplate{headline}{}
\setbeamertemplate{blocks}[rounded][shadow=false]
\definecolor{mmpageWhite}{HTML}{FCFBF7}
\definecolor{mmtanSoft}{HTML}{F7F5EF}
\definecolor{mmgreenDeep}{HTML}{244B2C}
\definecolor{mmgreenLeaf}{HTML}{5D8A56}
\definecolor{mmgreenSoft}{HTML}{E4EFD9}
\definecolor{mmtext}{HTML}{163221}
\definecolor{mmwarn}{HTML}{8F2F36}
\definecolor{mmwarnSoft}{HTML}{F8E8EA}
\setbeamercolor{background canvas}{bg=mmpageWhite}
\setbeamercolor{normal text}{fg=mmtext,bg=mmpageWhite}
\setbeamercolor{block title}{fg=mmgreenDeep,bg=mmgreenSoft}
\setbeamercolor{block body}{fg=mmtext,bg=mmtanSoft}
\setbeamercolor{block title example}{fg=mmgreenDeep,bg=mmgreenSoft}
\setbeamercolor{block body example}{fg=mmtext,bg=white}
\setbeamercolor{block title alerted}{fg=mmwarn,bg=mmwarnSoft}
\setbeamercolor{block body alerted}{fg=mmtext,bg=white}
\setbeamercolor{itemize item}{fg=mmgreenDeep}
\newtcolorbox{MMCard}[2][]{enhanced,breakable=false,colback=#2,colframe=black,boxrule=1.5pt,arc=8pt,left=5pt,right=5pt,top=4pt,bottom=4pt,before skip=0pt,after skip=0pt,#1}
\newcommand{\KoruIcon}{\raisebox{-0.2em}{\includegraphics[height=1.05em]{../../../manamaths/SITE/header-logo.png}}}
\newcommand{\NotesTitle}[1]{{\colorbox{mmtanSoft}{\parbox{0.975\linewidth}{\vspace{0.14em}\hspace{0.25em}{\LARGE\bfseries\textcolor{mmgreenDeep}{#1}\hfill\KoruIcon}\vspace{0.14em}}}\par\vspace{0.18em}{\color{mmgreenLeaf}\rule{\linewidth}{2.0pt}}\par\vspace{0.12em}}}
\newcommand{\SectionTitle}[1]{{\large\bfseries\textcolor{mmgreenDeep}{#1}}\par\vspace{0.04em}}
\newcommand{\GreenDivider}{\par\vspace{0.01em}{\color{mmgreenLeaf}\rule{\linewidth}{2.4pt}}\par\vspace{0.03em}}
\newcommand{\KeyIdea}[1]{\begin{MMCard}[title={\textbf{Key idea}},colbacktitle=mmgreenSoft,coltitle=mmgreenDeep]{mmtanSoft}\small #1\end{MMCard}}
\newcommand{\WorkedExample}[2]{\begin{MMCard}[title={\textbf{#1}},colbacktitle=mmgreenSoft,coltitle=mmgreenDeep]{white}\small #2\end{MMCard}}
\newcommand{\CommonMistake}[1]{\begin{MMCard}[title={\textbf{Common mistake}},colbacktitle=mmwarnSoft,coltitle=mmwarn]{white}\small #1\end{MMCard}}
\newcommand{\TryThese}[1]{\begin{MMCard}[title={\textbf{Try these}},colbacktitle=mmgreenSoft,coltitle=mmgreenDeep]{mmtanSoft}\small #1\end{MMCard}}
\newcommand{\StepLine}[2]{\textbf{#1.} #2\par\vspace{0.03em}}
\usepackage{tikz}
\setbeamertemplate{background}{\begin{tikzpicture}[remember picture,overlay]\draw[black,line width=1.5pt,rounded corners=10pt]([xshift=0.45cm,yshift=-0.45cm]current page.north west) rectangle ([xshift=-0.45cm,yshift=0.45cm]current page.south east);\end{tikzpicture}}
\begin{document}
\begin{frame}[t,shrink=10]
\NotesTitle{Notes \& Steps}
\footnotesize
\begin{columns}[T,onlytextwidth]
\begin{column}{0.48\textwidth}
\KeyIdea{Substitution means replacing a variable with its value, then simplifying carefully. If \(x=3\), then every \(x\) is replaced by \(3\).}
\SectionTitle{Steps}
\StepLine{1}{Write the value of the variable clearly.}
\StepLine{2}{Replace each variable with that value.}
\StepLine{3}{Use brackets if needed, especially with negatives.}
\StepLine{4}{Work out the expression using GEMA.}
\end{column}
\begin{column}{0.48\textwidth}
\SectionTitle{Key facts}
\begin{itemize}
  \item If \(x=4\), then \(2x=8\).
  \item If \(a=5\), then \(a^2=25\).
  \item If \(p=3\) and \(q=2\), then \(p+2q=7\).
  \item Brackets matter: if \(x=-2\), then \(x^2=4\) but \(-x^2=-4\).
\end{itemize}
\end{column}
\end{columns}
\GreenDivider
\begin{columns}[b,onlytextwidth]
\begin{column}{0.48\textwidth}
\CommonMistake{Forgetting brackets with negatives. If \(x=-3\), then \(x^2=(-3)^2=9\), not \(-9\).}
\end{column}
\begin{column}{0.48\textwidth}
\TryThese{\begin{enumerate}
  \item If \(x=3\), find \(x+4\).
  \item If \(a=5\), find \(2a\).
  \item If \(m=2,n=4\), find \(m+n\).
\end{enumerate}}
\end{column}
\end{columns}
\end{frame}
\begin{frame}[t,shrink=12]
\NotesTitle{Notes \& Steps}
\begin{columns}[T,onlytextwidth]
\begin{column}{0.48\textwidth}
\WorkedExample{Example 1: one variable}{If \(x=4\), find \(2x+3\).\[2x+3=2(4)+3=8+3=11\]}
\end{column}
\begin{column}{0.48\textwidth}
\WorkedExample{Example 2: two variables}{If \(p=5\) and \(q=2\), find \(p+2q\).\[p+2q=5+2(2)=5+4=9\]}
\end{column}
\end{columns}
\GreenDivider
\begin{columns}[T,onlytextwidth]
\begin{column}{0.48\textwidth}
\WorkedExample{Example 3: powers}{If \(a=3\), find \(a^2+1\).\[a^2+1=3^2+1=9+1=10\]}
\end{column}
\begin{column}{0.48\textwidth}
\WorkedExample{Example 4: negative value}{If \(x=-2\), find \(x^2+3\).\[x^2+3=(-2)^2+3=4+3=7\]}
\end{column}
\end{columns}
\end{frame}
\end{document}
